% CPE333 Operating Systems - PS7 Concurrency and Threads
% Standalone report source. Compile with pdfLaTeX.
\documentclass[a4paper,11pt]{article}
\usepackage[utf8]{inputenc}
\usepackage[T1]{fontenc}
\usepackage{lmodern}
\usepackage[a4paper,margin=2.2cm]{geometry}
\usepackage{graphicx}
\usepackage{float}
\usepackage{caption}
\usepackage{booktabs}
\usepackage{hyperref}
\usepackage{xcolor}
\usepackage{listings}
\hypersetup{colorlinks=true,linkcolor=blue!60!black,citecolor=blue!60!black,urlcolor=blue!60!black}
\lstset{
  basicstyle=\ttfamily\small,
  backgroundcolor=\color{gray!8},
  frame=single,
  breaklines=true,
  columns=fullflexible,
  keepspaces=true,
  showstringspaces=false,
  captionpos=b
}
\title{CPE333 Operating Systems\\Problem Session 7: Concurrency and Thread}
\author{Member 1: Name Thanaboon Tikaew 67070501021 \\
Member 2: Chanon Lhumsa-ard 67070501059 \\
Member 3: Theeraphat Jaingam 67070501063 \\
Member 4: Siriwan Yindeepot 67070501075}
\date{September 29, 2026}

\begin{document}
\maketitle

\begin{center}
\fbox{\parbox{0.92\linewidth}{\textbf{Environment:} Ubuntu 24.04 on WSL2, GCC 13.3.0, POSIX threads. The single source file \texttt{ps7.c} was compiled and run in WSL. Task 2 and Task 3 results below are five independent runs captured on September 29, 2026.}}
\end{center}

\section{Program design}
\subsection{Example and shared variable}
The program models available seats in a reservation system. The shared variable \texttt{available\_seats} starts at 100,000. A reservation decreases the value by one, and a cancellation increases it by one. Each activity performs 100,000 operations, so without lost updates the expected final value is 100,000. All three tasks are implemented in the single source file \texttt{ps7.c}; a command-line argument selects which task to run.

From the \texttt{PS7} directory in Ubuntu, compile the program once and then run each task using its number. Each run starts a fresh process with the counter reset to 100,000.
\begin{lstlisting}
cd ~/CPE333/PS7
gcc -o ps7 ps7.c -pthread
./ps7 1   # Task 1
./ps7 2   # Task 2
./ps7 3   # Task 3
\end{lstlisting}

\section{Code and Execution Results}
The program remains in one source file, \texttt{ps7.c}. The task-specific code excerpts and their corresponding results are shown separately below.

\subsection{Task 1: one-thread baseline}
One worker performs the reservation loop followed by the cancellation loop. The main thread waits for the worker to finish.
\begin{lstlisting}[basicstyle=\ttfamily\scriptsize,caption={Task 1 code excerpt}]
static void *single_thread_baseline(void *unused)
{
    (void)unused;

    for (long i = 0; i < OPERATIONS_PER_THREAD; ++i) {
        available_seats = available_seats - 1;
    }
    for (long i = 0; i < OPERATIONS_PER_THREAD; ++i) {
        available_seats = available_seats + 1;
    }

    return NULL;
}
\end{lstlisting}
\textbf{Result (fresh run):}
\begin{table}[H]
\centering
\caption{Task 1 output from \texttt{./ps7 1}.}
\begin{tabular}{lr}
\toprule
Metric & Available seats \\
\midrule
Initial count & 100,000 \\
Expected final count & 100,000 \\
Observed final count & 100,000 \\
\bottomrule
\end{tabular}
\end{table}
\textbf{Summary of values:} The final count matched the expected 100,000 seats.

\subsection{Task 2: concurrent updates without deliberate yielding}
The shared worker routine loads, updates, and stores the counter. Task 2 passes a zero context-switch flag, so the worker does not call \texttt{sched\_yield()}.
\begin{lstlisting}[basicstyle=\ttfamily\scriptsize,caption={Task 2 code excerpt}]
static void *update_seats(void *arg)
{
    const worker_args_t *worker = (const worker_args_t *)arg;

    for (long i = 0; i < OPERATIONS_PER_THREAD; ++i) {
        long local_seats = available_seats;
        if (worker->force_context_switch) {
            sched_yield();
        }
        local_seats += worker->change;
        available_seats = local_seats;
    }

    return NULL;
}

/* Task 2 selection in main(): */
case 2:
    return run_race_task(0, 2);
\end{lstlisting}
\textbf{Result} (final available seats from five independent runs):
\begin{table}[H]
\centering
\caption{Task 2 final available seats in five runs.}
\begin{tabular}{cr}
\toprule
Run & Final available seats \\
\midrule
1 & 54,358 \\
2 & 116,800 \\
3 & 0 \\
4 & 75,290 \\
5 & 199,968 \\
\bottomrule
\end{tabular}
\end{table}
\textbf{Summary of values:} The observed final count ranged from 0 to 199,968.

\clearpage
\subsection{Task 3: concurrent updates with \texttt{sched\_yield()}}
Task 3 reuses the worker routine shown for Task 2, but selects the context-switch flag as one. This enables the \texttt{sched\_yield()} call between the load and store.
\begin{lstlisting}[basicstyle=\ttfamily\scriptsize,caption={Task 3 code excerpt}]
/* Worker flags created by run_race_task(1, 3): */
worker_args_t workers[2] = {
    {-1, 1},
    {+1, 1},
};

/* Task 3 selection in main(): */
case 3:
    return run_race_task(1, 3);
\end{lstlisting}
\textbf{Result} (final available seats from five independent runs):
\begin{table}[H]
\centering
\caption{Task 3 final available seats in five runs with \texttt{sched\_yield()}.}
\begin{tabular}{cr}
\toprule
Run & Final available seats \\
\midrule
1 & 52,845 \\
2 & 79,635 \\
3 & 110,786 \\
4 & 0 \\
5 & 199,928 \\
\bottomrule
\end{tabular}
\end{table}
\textbf{Summary of values:} The observed final count ranged from 0 to 199,928.

\section{Discussion}
Task 1 gives the expected result because one worker updates the shared variable sequentially. In Tasks 2 and 3, each update is split into a load, a calculation, and a store. For example, if the counter is 100 and both workers read it, the reservation worker calculates 99 while the cancellation worker calculates 101. If the cancellation worker stores last, its value overwrites the reservation update and the final count is 101 instead of 100. Reversing the stores can leave 99. This lost update is the race condition.

Task 2 varies even without an explicit yield because the workers can run concurrently and their operations can interleave naturally. In these five runs, the final values ranged from 0 to 199,968. Task 3 places \texttt{sched\_yield()} between the load and store, giving another runnable worker an opportunity to run. The five observed values ranged from 0 to 199,928. These results show scheduling-dependent behavior; they do not establish a guaranteed direction or final value. \texttt{sched\_yield()} requests a scheduling opportunity but does not guarantee a context switch.

\textbf{Technical note:} The shared variable is declared \texttt{volatile} so the demonstration keeps explicit loads and stores. \texttt{volatile} does not make an update atomic or synchronize the workers. Unsynchronized conflicting accesses are a C data race, so the language does not guarantee a particular result. A real reservation system should protect the shared counter with a mutex or an appropriate atomic operation.
\end{document}
